\chapterhead{2}{Sieves and Magnets}{Using a difference between the parts}\label{ch2}

Sorting by hand works for a handful of rice. For sand or for a thousand grains it is too slow. Instead we look for something that is \emph{different} about the parts and use a tool that notices that difference.

\begin{center}
\sffamily\small\renewcommand{\arraystretch}{1.5}
\begin{tabular}{@{}p{3.0cm}p{3.6cm}p{3.6cm}@{}}
\toprule
\textbf{Tool} & \textbf{It notices a difference in} & \textbf{It works when}\\
\midrule
A sieve & size & one part is small enough to fall through\\
A magnet & steel or iron & one part is steel and the rest is not\\
Your hands & looks and shape & there are not many parts\\
\bottomrule
\end{tabular}
\end{center}

\begin{tryit}{Sieve the sand}
\textbf{You need:} a kitchen scale, a sieve or tea strainer, a bowl, 150~g of sand and 60~g of lentils.
\begin{enumerate}[leftmargin=1.6em]
\item Mix the sand and the lentils. Weigh the whole mixture.
\item Hold the sieve over a bowl. Pour the mixture in and shake gently.
\item Weigh what fell through and what stayed on the sieve.
\end{enumerate}
\par\medskip
\renewcommand{\arraystretch}{1.6}
\begin{tabular}{|l|c|c|c|c|}
\hline
 & \textbf{Whole} & \textbf{Fell through} & \textbf{Stayed on sieve} & \textbf{Parts added}\\ \hline
\textbf{Weight (g)} & & & & \\ \hline
\end{tabular}
\end{tryit}

\begin{worked}{Practice numbers, not from a real test}
Whole mixture: 210~g. Fell through: 148~g of sand. Stayed on: 60~g of lentils.\\
$148+60=208$. The gap is $210-208=2$~g. Some sand stuck between the lentils, and a pinch dropped on the table.
\end{worked}

\begin{tryit}{Use a magnet}
Mix 100~g of rice with 20~g of paper clips. Hold a magnet over the mixture, or pull it through. Weigh what it picks up and what is left.\par\medskip
\renewcommand{\arraystretch}{1.6}
\begin{tabular}{|l|c|c|c|}
\hline
 & \textbf{Whole} & \textbf{On the magnet} & \textbf{Left behind}\\ \hline
\textbf{Weight (g)} & & & \\ \hline
\end{tabular}
\end{tryit}
\adult{Small magnets can harm if swallowed. Keep them away from young children and pets. Some paper clips are plated and are not magnetic, so test one first.}

\begin{worked}{Practice numbers}
Whole: 120~g. The magnet picked up 18~g. The rest weighed 102~g. $18+102=120$.\\
Two grams of clips stayed in the rice. They were a kind of clip that is not magnetic. The total still matched. Only the \emph{split} was wrong.
\end{worked}

\begin{multix}[hmm]\multixbot[hmm]\tcblower
A sieve has holes, so it can separate salt from sand. Table salt and sand grains are about the same size, but holes are holes. Will it work?
\end{multix}

\trust{We weighed the whole and the parts after sieving and after the magnet, and we checked what stayed behind in each tool.}
{The totals matched within a couple of grams. The magnet check also showed us a hidden clip that the magnet could not pick up.}
{That a tool separates a mixture completely. A sieve or a magnet is only as good as the difference it relies on.}

\begin{keepbox}
We used a sieve and a magnet. \textbf{We kept} \wline[3.4cm] \textbf{and we set aside} \wline[3.4cm].
\end{keepbox}

\clearpage

\begin{puzzle}{Choose the tool}
Write S for sieve, M for magnet, H for hand or X if none of these will work.\\
(a) sand and pebbles \wline[1cm] \quad (b) rice and steel screws \wline[1cm] \quad (c) salt and sand \wline[1cm]\\
(d) sugar dissolved in tea \wline[1cm] \quad (e) flour and big stones \wline[1cm]
\end{puzzle}

\begin{puzzle}{Add up}
A sieve test started with 210~g. Sand through the sieve weighed 146~g and lentils on it 60~g. Total \wline[1.6cm]~g. Gap \wline[1.6cm]~g.
\end{puzzle}

\begin{puzzle}{Magnet test}
A magnet picked up 15~g from 90~g of rice and clips. The rest weighed 75~g. Add them. \wline[1.6cm]\\
Does it match? \wline[2cm]
\end{puzzle}

\begin{puzzle}{Hidden clip}
After the magnet, a student found a clip still in the rice. Which kind of clip is it? \wline[4cm] Where did the gap come from? \wline[4cm]
\end{puzzle}

\begin{puzzle}{Unfair test}
Ana sieved damp sand and Ben sieved dry sand. Ana's sieve clogged. They say the sieve is bad. What is wrong with their test? \wline[6cm]
\end{puzzle}

\clearpage

\begin{puzzle}{Order the steps}
Number the steps for the sand and lentils from 1 to 4.\\
\wline[1cm] Weigh what fell through. \quad \wline[1cm] Weigh the whole mixture.\\
\wline[1cm] Add the two parts. \quad \wline[1cm] Shake the sieve.
\end{puzzle}

\begin{puzzle}{Which property?}
For each tool, write the property it uses.\\
Sieve: \wline[3cm] \quad Magnet: \wline[3cm]
\end{puzzle}

\begin{puzzle}{Does the order matter?}
Rice, clips and sand are mixed. Would you use the magnet first or the sieve first? Say why either order could work. \wline[6cm]
\end{puzzle}

\begin{puzzle}{Draw the circles}
A mixture of 60, 25 and 10~g is sieved into 60 and 35. Fill in the empty circle.
\begin{center}
\begin{tikzpicture}
  \node[bigcir,minimum size=2.4cm] (a) at (0,0) {}; \node[font=\sffamily\scriptsize] at (0,0) {mixture 95};
  \node[cir] (b) at (4,0.8) {60}; \node[cir,dashed] (c) at (4,-0.8) {};
  \draw[popA] (a) -- (b); \draw[popA] (a) -- (c);
\end{tikzpicture}
\end{center}
\end{puzzle}

\begin{puzzle}{Design a sorter}
Choose a mixture at home and say which difference between its parts you could use. \wline[6cm]
\end{puzzle}
\clearpage

\thispagestyle{fancy}
\subhead{Sieve record}
Try three different mixtures through the same sieve. Weigh before and after, and record what stayed on and what fell through.\par\medskip
\renewcommand{\arraystretch}{3.0}
\begin{tabularx}{\linewidth}{|>{\sffamily\small\bfseries}p{2.6cm}|X|X|X|X|}
\hline
Mixture & Whole (g) & Fell through (g) & Stayed on (g) & Gap (g) \\ \hline
 & & & & \\ \hline
 & & & & \\ \hline
 & & & & \\ \hline
\end{tabularx}

\vspace{8pt}
Which mixture split most cleanly? \wline[7cm]\\[10pt]
Which left the biggest gap, and where might the gap be? \wline[5cm]\\[10pt]
Sketch your sieve and mark where each part went.\par\medskip
\begin{tikzpicture}\draw[ink!70,line width=1pt,rounded corners=3pt] (0,0) rectangle (11.4,5.2);\end{tikzpicture}
\clearpage
